\documentclass{article}

% ready for submission
% \usepackage[final]{aaj2026}
\bibliographystyle{plainnat}
\usepackage[numbers]{natbib}
\usepackage[utf8]{inputenc} % allow utf-8 input
\usepackage[T1]{fontenc}    % use 8-bit T1 fonts
\usepackage{hyperref}       % hyperlinks
\usepackage{url}            % simple URL typesetting
\usepackage{booktabs}       % professional-quality tables
\usepackage{amsfonts}       % blackboard math symbols
\usepackage{nicefrac}       % compact symbols for 1/2, etc.
\usepackage{microtype}      % microtypography
\usepackage{xcolor}         % colors
\usepackage{amsmath}

\title{Fast Monte Carlo Pricing of Asian Options under Rough Fractional Stochastic Volatility}

% The \author macro works with any number of authors. There are two commands
% used to separate the names and addresses of multiple authors: \And and \AND.
%
% Using \And between authors leaves it to LaTeX to determine where to break the
% lines. Using \AND forces a line break at that point. So, if LaTeX puts 3 of 4
% authors names on the first line, and the last on the second line, try using
% \AND instead of \And before the third author name.


\author{%
  Tiago Flora \\
  Department of Mathematics\\
  University of Chicago \\
  \texttt{tflora@uchicago.edu} \\
  % examples of more authors
}


\begin{document}


\maketitle

  \begin{abstract}  We study Monte Carlo pricing of a fixed-strike arithmetic Asian call option under the
  Rough Fractional Stochastic Volatility (RFSV) model \citep{gatheral2018volatility},
  \end{abstract}


  \section{Motivation}

  \subsection{Asian options and rough volatility}
  In this project, I am interested in pricing fixed-rate, arithmetic-average Asian call options, which have payoff $C(T) = \max (A(0,T)-K,0)$, where $A(0,T) = \frac 1 T \int^T_0 S_tdt$, and $S_t$ is the price of the underlying security with respect to time. Importantly, I assume that $\log \sigma_t$, or the log-volatility of $S_t$, follows a fractional Brownian motion with Hurst parameter 0.1, following the famous result from Gatheral, Jaisson \& Rosenbaum in \cite{gatheral2018volatility}. I also evaluate simulations with other values of $H$.

  If we formulate the problem as a function of discrete prices, such that $A(0,T) = \frac 1 T \sum^T_j S_{t_j}$, then even under Black-Scholes dynamics each $S_{t_j}$ is a lognormal random variable, and the sum of lognormal random variables is not lognormal. The probability density function for the sum of correlated lognormal variables has no known exact, closed-form analytical representation. Our case is worse: under rough fractional stochastic volatility (RFSV), $S_t$ is no longer lognormal, and we cannot use approximations like L\'{e}vy's moment-matching approximations \cite{levy1992pricing}.

 Under RFSV, we have $\log \sigma_t = \mu + \nu W^H_t$. Because the volatility paths are (negatively) correlated and rough, calculating the exact joint expectation $E[S_{t_i} S_{t_j}]$ for the second moment becomes complex and solving the moment-matching equations practically impossible: it would require computing the expectation of the integral of the volatility path, which always depends on the entire continuous history of the volatility up to that point. The ``vol-of-vol" ($\nu$) parameter also introduces heavy skew and fat tails into the price distribution that a simple 2-moment lognormal approximation cannot capture.

 In these circumstances, Monte Carlo sampling becomes a viable option. It doesn't require a known probability density function; we can simulate $M$ paths of the RFSV process and calculate the payoff on each path. Our price estimate is simply given by:
 \begin{align}
    \lim_{M \to \infty} \frac{1}{M} \sum_{m=1}^M e^{-rT} \max\left(A^{(m)} - K, 0\right) = \mathbb{E}^{\mathbb{Q}}\left[e^{-rT} \max(A - K, 0)\right]
 \end{align}
    Where $\mathbb Q$ is the risk-neutral measure.

\subsection{The computational bottleneck of MC sampling for RFSV}
    
 Naïve MC sampling requires $O(N^3)$ setup work, for $N$ being the number of simulated time-steps until expiration: fBM is non-Markovian, so no forward recursion exists; every path requires drawing from $\mathcal{N}(0, C)$ where $C$ is a dense $N \times N$ positive-definite matrix. It then also requires $O(N^2)$ work per path, as we multiply the lower-triangular $L$ from the decomposition by a white noise vector to construct a correlated fBM path. For $N = 252$ trading days \& $M = 10{,}000$ paths, the $O(N^3)$ Cholesky factorization alone requires $N^3/3 \approx 5.3$ million FLOPs (the leading constant is $1/3$, versus $2/3$ for LU); the $O(N^2)$ per-path cost aggregates to a total $O(MN^2)$ workload, which dominates for $M$ large enough for reasonable accuracy. And since errors for MC sampling famously fall with $1/\sqrt M$, $M$ might have to be quite large depending on the spread size that matters to you.
 
This problem is relevant to any firm unit with Asian options in its portfolio; trading desks, for instance, often need to run parameter sweeps to build and update sensitivity surfaces, and manage real-time risk by running pricing engines with deviations to inputs. If one moves from inter- to intraday monitoring, $N$ grows by a large factor, and speeding up sampling can prove itself an unlikely life-saver.
  
\subsection{Approach and scope}

I implement and benchmark 3 algorithms, including the na\"ive implementation described in the previous section. 



To speed up both setup and MC path sampling, I tested 

% \begin{table}
%   \caption{Sample table title}
%   \label{sample-table}
%   \centering
%   \begin{tabular}{lll}
%     \toprule
%     \multicolumn{2}{c}{Part}                   \\
%     \cmidrule(r){1-2}
%     Name     & Description     & Size ($\mu$m) \\
%     \midrule
%     Dendrite & Input terminal  & $\sim$100     \\
%     Axon     & Output terminal & $\sim$10      \\
%     Soma     & Cell body       & up to $10^6$  \\
%     \bottomrule
%   \end{tabular}
% \end{table}


\medskip

\bibliography{biblio.bib}


%%%%%%%%%%%%%%%%%%%%%%%%%%%%%%%%%%%%%%%%%%%%%%%%%%%%%%%%%%%%

\appendix

\section{Appendix / supplemental material}


Optionally include supplemental material (complete proofs, additional experiments and plots) in appendix.
All such materials \textbf{SHOULD be included in the main submission.}

%%%%%%%%%%%%%%%%%%%%%%%%%%%%%%%%%%%%%%%%%%%%%%%%%%%%%%%%%%%%




\end{document}